\section{Some Useful Theorems}

\begin{lemma}\label{lem:gcd-ideals}
	Let $\O$ be the ring of integers of a number field and let $\a, \b$ be two ideals such that $\a + \b = \O$ (That is, thye are coprime ideals). 
	If $\a\b = \c^n$ for some ideal $\c$ and some integer $n \geq 1$, then $\a = \d^n$ and $\b = \e^n$ for some ideals $\d, \e$ such that $\d\e = \c$.
\end{lemma}

\begin{proof}
	This is direct consequence from unique factorisation of ideals into prime ideals.
\end{proof}

\begin{lemma}\label{lem:principal}
	Let $\a$ be an ideal of $\O$ and $m$ some decimal integers.
	If $gcd(h(\Qd), m) = 1$ and $\a^m$ is a principal ideal, then $\a$ is a principal ideal as well.
\end{lemma}

\begin{proof}
	We only need to recall that $h(\Qd)$ is the order of the ideal class group of $\Qd$, whose identity element is the class of principal ideals.

	By elementary theory of group theory we know that $\a^{h(\Qd)}$ must be idetity.
	Since $\a^m$ is also identity, we can write $1 = am + bh(\Qd)$ by Euclidean algorithm for some integers $a, b$ therefore $\a = \a^{am + bh(\Qd)} = (\a^m)^a (\a^{h(\Qd)})^b$ is also identity, which means $\a$ is principal.
\end{proof}

\begin{lemma}\label{lem:unit}
	The unit group of $\O_{\Qd}$ is $\{\pm 1\}$ for any even $d \geq 1$. 
\end{lemma}

\begin{proof}
	We know $u$ is a unit if and only if $N(u) = 1$ where $N$ is the norm function from $\Qd$ to $\Z$.

	We can write $u = a + b \smd$. Since $d$ is even, we can let $a, b \in \Z$.
	Then $N(u) = a^2 + db^2$, which equals $1$ if and only if $a = \pm 1$ and $b = 0$.
\end{proof}

\begin{conjecture}
	In equation \ref{eq:main}, $gcd(x, d) = 1$.
\end{conjecture}

\begin{lemma}\label{lem:automorphism}
	Let $G$ be a finite abelian group and $p$ a prime number such that $p \nmid |G|$, then the map $\phi: G \rightarrow G, \phi(g) = p\cdot g$ is an automorphism (bijective self isomorphism), where $p \cdot g$ denotes $\underbrace{g + g + g \cdots + g}_{p \text{ times}}$.
\end{lemma}

\begin{proof}
	It is without saying that $\phi$ is a homomorphism.

	To prove it is an automorphism, let us first prove the case when $G$ is cyclic.
	Say $G \cong \Z/n\Z$ for some integer $n$.
	Let $a, b \in G$ and assume that $\phi(a) = \phi(b)$, which is equivalent to 
	$$
		p a = p b \mod n \Leftrightarrow p(a - b) \equiv 0 \mod n
	$$
	Now, $p$ is a prime number and $p \nmid n$, therefore $a - b \equiv 0 \mod n$, which means $a = b$ and $\phi$ is injective.
	Since $G$ is finite, injectivity automatically implies surjectivity, therefore $\phi$ is an automorphism.

	In the general case, by the classification of finite abelian groups, we can write $G$ as a direct product of cyclic groups, i.e, $G \cong \Z/n_1\Z \times \Z/n_2\Z \times \cdots \times \Z/n_k\Z$ for some integers $n_1, n_2, \cdots, n_k$, where $p \nmid n_i$ for all $i$.
	Then we can apply the previous result to each cyclic component and conclude that $\phi$ is an automorphism on each component, therefore $\phi$ is an automorphism on the whole group $G$.
\end{proof}
